\section{Introduction}
\label{sec:intro}

An LLM-agent system in production is often thoroughly observable while governance
varies sharply by resource family.

The observability half is well served. Tracing frameworks capture spans,
tool calls, retrieved context, and token accounting; evaluation harnesses score
outputs offline; LLM-as-judge methods approximate human preference at a fraction
of the cost~\cite{zheng2023judging,liu2023geval}; and hosted platforms
render all of it in a dashboard~\cite{langsmith,langfuse,phoenix}. When an agent
produces a bad answer, a competent team will find out, and will be able to point
at the span where it happened.

What happens next is where the tooling stops. Someone opens a source file and
edits a prompt string. Perhaps they test it on a handful of examples they
assembled by hand; perhaps they compare against the previous version, or perhaps
they only check that the specific reported failure is fixed. They merge, deploy,
and move on. Six weeks later a different failure arrives, and nobody can
reconstruct what the prompt used to be, what evidence justified the change, who
approved it, or whether the change that fixed one class of input silently
degraded three others. The remediation path is manual, undocumented, and
irreproducible --- precisely the asymmetry Sculley et~al.\ described for
conventional ML systems~\cite{sculley2015debt}, displaced from model artifacts
onto prompts, tool definitions, agent skills, and retrieval corpora.

This gap is not an oversight in any individual tool; it is a consequence of how
the ecosystem is factored. Composition frameworks are built to answer \emph{how do
I build this}~\cite{langchain,llamaindex}. Program optimizers are built to answer
\emph{what should this prompt say}~\cite{khattab2024dspy,agrawal2025gepa,
opsahlong2024mipro}. Tracking and observability systems are built to answer
\emph{what happened} --- and, increasingly, more than that: \mlflow{}'s Prompt
Registry, LangSmith, and Langfuse now version prompts, expose a named pointer the
client resolves at call time, and roll back by re-pointing
it~\cite{mlflowpromptregistry,langsmithprompts,langfuseprompts}. The question
\emph{may this change go live, on what evidence, and how do I undo it} is
therefore answered today for one artifact type, by pointer discipline, with the
evidence sitting beside the decision rather than gating it.

What remains unanswered is the same question for everything else an agent depends
on --- its workflows, tool definitions, MCP server configurations, and retrieval
corpora --- and the question of what it would cost to answer it for all of them at
once. That is this paper's subject, and \secref{sec:compare} states the delta
against those platforms feature by feature rather than by category.

\para{Three words, used deliberately differently.} Because this paper distinguishes
things a looser treatment would conflate, we pin the vocabulary once and hold to it.
An \textbf{AI-agent resource} is something in the world an agent depends on at run
time: a prompt, a workflow, a tool, an MCP server, a corpus. An \textbf{artifact} is
the generic term for such a thing considered as a versioned object, used when the
contrast is with \emph{code}, which already has a release discipline. A
\textbf{\gasset} is \sys's typed record \emph{for} one --- the abstraction this paper
introduces --- and only the third term is technical. Where the paper says
``artifact'' it is never referring to an \mlflow{} artifact-store object;
\appref{app:config} notes that collision explicitly.

\para{This paper.} We present \sysbare{}, a layered control plane for the
lifecycle of AI-agent resources. \sys{}'s organizing abstraction is the
\emph{\gasset}: a typed record carrying a schema-validated definition, version
history, an authority model, a trace and audit trail, and --- where the asset's
family supports one --- a \livetarget{} that client code resolves at call time,
so a release changes behaviour without a client deployment. Around that
abstraction \sys arranges \statLayersW{} layers and \statModesW{} reusable lifecycle
modes, and composes them into a \chainname{} that carries a flagged production
trace to a measured candidate, an enforced evaluation gate, an explicit human
decision, and an audited release that records the outgoing target so rollback is a
restore rather than a reconstruction.

\para{The central claim, stated narrowly.} It would be easy to claim that a
layered control plane makes governance uniform across artifact types. That claim
would be false, and the interesting content of this paper is why. Our actual
claim is about what layering does and does not buy:

\begin{quote}\itshape
Adjacency in a layered architecture confers capability \emph{availability}, not
capability \emph{inheritance}. A family placed in the asset layer obtains
lifecycle behaviour and governance enforcement only by explicitly wiring them.
The shared base contract is identity, history, authority, and audit; release,
rollback, evaluation, and evidence are capability-specific obligations. \sys
declares those obligations per family in a checked registry, but does not yet tie a
declaration to the handler that discharges it (\secref{sec:arch:claim}).
\end{quote}

We defend the per-family capability factoring while treating the remaining gap ---
declarations that are checked against each other but not against the handlers
behind them --- as technical debt (\secref{sec:arch}). A prompt
declares live-target release and rollback; a test set is evidence and declares no
live target; a judge is evaluatable but not releasable. We therefore give the
guarantee surface in full, including the rows where it is weaker than a reader
might assume (\tabref{tab:families}).

\para{Contributions.}
\begin{enumerate}
\item \textbf{A typed, per-family governance architecture.} A
\statLayersW{}-layer factoring in which the unit of governance is an asset rather
than a model or trace, plus an explicit capability/guarantee matrix that separates
the shared base contract from family-specific release, rollback, and evidence
semantics (\secref{sec:overview}, \secref{sec:arch}).

\item \textbf{A \chainname{} distinguished from its implementations.} A
seven-term reference shape --- \cSignal{}, \cEvidence{}, \cCandidate{},
\cMeasure{}, \cDecision{}, \cRelease{}, \cTrace{} --- with an explicit account of
how the concrete six-stage prompt path realizes it non-bijectively, and of what
durable residue each term leaves so that a release is reconstructable from records
(\secref{sec:overview}, \figref{fig:chain}).

\item \textbf{A gate taxonomy that survives contact with operators.} The
separation of an \emph{enforced} \advgate{} from an \emph{advisory} per-version
release verdict, and the argument that putting the enforced gate on candidate
advancement rather than on release is what keeps operators from routing around the
system (\secref{sec:capabilities}, D3 in \tabref{tab:decisions}).

\item \textbf{A stated per-loop guarantee surface for a broker-free execution
model.} Claim-and-lease columns on relational rows are a known
pattern~\cite{kleppmann2017ddia} and collapsing the worker tier into the API
process is an engineering choice rather than a research result; we claim neither as
novel. What we do contribute is the discipline of stating what each of the
\statLoopsW{} loops actually guarantees --- they differ, and an earlier draft of
this paper wrongly described them with one number (\tabref{tab:loops}) --- together
with the honest converse that in-memory rate limits become per-process and must be
sized for the worker count (\secref{sec:exec}).

\item \textbf{Precise state ownership, including its limits.} Three stores, three
owners, and the two dual-write boundaries no transaction can span. Naming those
boundaries is what allows a precise audit guarantee to be stated at all
(\secref{sec:components}, \figref{fig:dataflow}).

\item \textbf{A comparison that includes the losses.} \sys against tracking,
composition, multi-agent, and optimizer systems along \statCompareDims{}
dimensions, with a cost column for every row including \sys's own
(\secref{sec:compare}).

\item \textbf{A measurement protocol, offered as an artifact rather than as a
result.} \appref{app:protocol} specifies the questions, the baselines, the
fault-injection matrix, and the analysis in enough detail to be executed by someone
else. We want to be exact about its standing: an unpopulated table is not a
contribution, and we do not present it as one. The protocol is a usable artifact;
the results would be the contribution, and they are absent
(\secref{sec:eval}).
\end{enumerate}

\para{What this paper does not claim.} \sys is not a composition or multi-agent
orchestration framework, and does not compete with one; agents are composed
elsewhere and governed here. It does not rebuild tracking: \mlflow{}'s
experiments, traces, assessments, prompt registry, artifact store, and evaluation
entry points are dependencies, not
reimplementations~\cite{mlflowgenai,mlflowpromptregistry}. It ships a single environment with no
\code{dev}\arrowto\code{staging}\arrowto\code{prod} ladder. Its local tool and MCP
execution is \emph{containment}, not isolation, and we are careful throughout
about which word applies (\secref{sec:limitations}). Most importantly, the
quantitative evaluation in \secref{sec:eval} is specified but not yet executed.
We state that plainly and mark each unpopulated result, because a systems paper
that decorates an unrun experiment with plausible numbers is worse than one that
reports an empty table.
